\documentclass[11pt,a4paper]{article}

\usepackage[utf8]{inputenc}
\usepackage[T1]{fontenc}
\usepackage{lmodern}
\usepackage[margin=2.4cm]{geometry}
\usepackage{amsmath}
\usepackage{amssymb}
\usepackage{xcolor}
\usepackage{enumitem}
\usepackage{comment}
\usepackage{booktabs}
\PassOptionsToPackage{hyphens}{url}
\usepackage[colorlinks=true,linkcolor=black,urlcolor=blue,citecolor=black]{hyperref}

\setlength{\parindent}{0pt}
\setlength{\parskip}{0.6em}
\emergencystretch=3em
\sloppy

% ---------------------------------------------------------------------------------------
% Internal notes.
%
% These record what still has to happen before a paragraph is truthful, and they must not
% reach the reviewers. They are excluded by default; switch the line below to
% \includecomment{internal} only while drafting.
% ---------------------------------------------------------------------------------------
\excludecomment{internal}
\newcommand{\internalhead}{\textcolor{red}{\textbf{[Internal note: remove before
submission]}}\par\nobreak}

% Quoted reviewer comment.
\newenvironment{reviewercomment}
  {\begin{quote}\itshape\small}
  {\end{quote}}

% One reviewer point: \point{label}{quoted comment}
\newcommand{\point}[2]{%
  \subsection*{#1}%
  \begin{reviewercomment}#2\end{reviewercomment}}

\title{Response to reviewers}
\author{}
\date{}

\begin{document}

\begin{center}
{\LARGE\bfseries Response to reviewers}\par\medskip
\end{center}

\textbf{Manuscript:} \emph{Omnisolver: An extensible interface to Ising spin--glass and QUBO
solvers -- a numerically stabilized, distributed GPU exhaustive-search plugin} (SoftwareX,
Original Software Publication update)

\textbf{Authors:} K.~Ja\l{}owiecki, J.~Paw\l{}owski, B.~Gardas, \L{}.~Pawela

\hrule

\begin{internal}
\internalhead
All six experiments (E1--E6) have been performed on $8\times$ H100 across two nodes, and their
results are in the manuscript. What remains is purely administrative: tag release
\texttt{0.0.6}, deploy the plugin documentation, and make the data repository publicly
reachable. Statuses are tracked in \texttt{REVIEW\_ANALYSIS.md}~\S0. Set
\texttt{\textbackslash excludecomment\{internal\}} in the preamble to hide every note like
this one.
\end{internal}

\hrule

We thank all three reviewers for a careful and constructive reading. The revision reframes the
contribution, adds five tables, a section on reproducing the results, several new paragraphs
and a revised precision analysis. We are particularly grateful to Reviewer~3 for pointing out
that the plugin was already announced in the original publication. Acting on that comment led
us to rewrite the contribution claim from scratch.

Two numbers from the submitted version have changed:

\begin{enumerate}[leftmargin=*]
\item \textbf{Speedup.} The parallel-efficiency table requested by Reviewer~1 gives $43\%$
efficiency at $N=44$, so the ideal $8\times$ speedup is now dated from $N=52$: the two samplers
cross between $N=40$ and $N=42$, efficiency exceeds $90\%$ from $N=48$, and saturates from
$N=52$ onwards.
\item \textbf{Precision.} The verification tables cover $N=40,\dots,54$ (single-GPU) and $N=60$
(distributed), and that paragraph is now written around those instances. With the CPU heuristic
described under R2.1 it covers all twenty stored runs, $N=38,\dots,60$, a wider range than
before.
\end{enumerate}

A \texttt{latexdiff} against the submitted version is provided.

\section*{Reviewer 1}

\point{R1.1}{The paper measures up to 8 H100 GPUs, but claims about scaling to much larger
allocations are back-of-the-envelope projections. The authors also note that controller merge
bandwidth, Ray overhead, and kernel launch latency may limit scaling.}

\textbf{We agree with this assessment, and we have reduced both the size of the projection
 and the emphasis placed on it.}
Three changes:

\begin{itemize}[leftmargin=*]
\item The projection target is now $2^{14}=16{,}384$ GPUs rather than $2^{16}$, which is the
scale we can defend as within reach of current leadership-class systems. The revised paragraph
gives $N=60$ in $\approx\!2.2$~min and $N=64$ in $\approx\!35$~min at $2^{14}$ GPUs, with a
one-month budget reaching $N \lesssim 74$. $2^{16}$ now appears only as an explicitly
hypothetical bound, with the observation that quadrupling the resources buys two variables.
\item The projection is confined to one labeled paragraph (\emph{Scaling outlook on larger
allocations}) and no longer appears in the abstract or the conclusions. It is now preceded by a
new paragraph, \emph{Controller cost and the limits of the decomposition}, which measures the
local controller components and identifies plausible limits (see R2.2), and it closes by
pointing back to that caveat.
\end{itemize}

The projection is now backed by measurement. Table~5 reports the
controller cost up to $2^{16}$ subproblems, $7.4$~s merging directly and $3.8$~s
hierarchically, of which the merge \emph{arithmetic} is only $2.1$~s. That measurement runs on
one machine, so it bounds the local component from below rather than estimating the total, as
the text and caption state. Table~4 measures every GPU count and node layout at fixed problem
size, giving $2.00\times$ on two GPUs, $3.99\times$ on four and $7.86\times$ on eight, i.e.\
$98.2\%$ parallel efficiency, with weak scaling flat to within $0.2\%$. The whole series was
measured twice, on separate days.

What remains extrapolated is behavior beyond eight devices and beyond one node boundary,
which is what \emph{Limitations and intended scope} now states.

\point{R1.2}{The paper positions the plugin as a ground-truth oracle and briefly checks results
against SBM, but it does not provide a broad benchmark against multiple exact or heuristic
solvers across different instance families.}

\textbf{We agree, and have added a second, algorithmically distinct heuristic, evaluated
against certified optima on four instance families.}

Since the plugin enumerates every assignment, its own runtime is essentially independent of
coefficient family. The revision restricts that claim to this plugin, since branch-and-bound
and structure-aware methods may prune the search and can depend strongly on the family.

The new paragraph \emph{Instance families and targeted multi-solver comparison} and Table~6
report dSB and \texttt{dwave.samplers.SimulatedAnnealingSampler} on twenty certified $N=44$
instances, five each with uniform, bimodal, Gaussian and sparse couplings, under the
method-specific budgets documented there. Both reached a certified optimum on all twenty. The
few states that differ from the stored certificate are degenerate ground states at the same
integer energy, and independent \texttt{float64} evaluation places every selected energy within
$3.2\times10^{-13}$ of the certificate. The complete records ship as
\path{benchmarks/results/multi_solver/run_20260816_125511.json}.

We also added the complementary case, a family on which the heuristic \emph{fails}. On the
Wishart ensemble (Hamze et al., Phys.\ Rev.\ E \textbf{101}, 052102 (2020)) at $N=40$, dSB
under the same budgets misses the certified optimum on almost every instance once the
ruggedness parameter reaches $\alpha\leq0.3$, terminating up to $0.6\%$ above optima known
only because brute force found them, while certification takes seconds. A paired planted arm
with a closed-form $E_0$ gives indistinguishable counts, so the difficulty is not a planting
artifact, and its certificates match the analytic $E_0$ to $7\times10^{-15}$. Fig.~2(b) puts
certification two to three orders of magnitude ahead in time-to-solution at the two hardest
points. All of this is the new Fig.~2 and the paragraph \emph{A rugged family where
certification is decisive}.

For implementation correctness we retain the comparison with \texttt{dimod.ExactSolver} at
CPU-enumerable sizes and independent \texttt{float64} recomputation at benchmark scale.

\point{R1.3}{The paper claims near-ideal speedup, but does not explicitly calculate or report
parallel efficiency.}

\textbf{We have added the requested parallel-efficiency table.}

Table~3 of the revision reports the speedup $S=t_1/t_8$ and the efficiency $E=S/8$ size by
size, computed from the raw result files, since both samplers ran on identical instances. It
dates the full $\approx\!8\times$ speedup from $N=52$ and not from $N \gtrsim 44$: efficiency
is $43\%$ at $N=44$, exceeds $90\%$ from $N=48$ and saturates from $N=52$ on. The abstract,
\emph{Architecture}, \emph{Performance} and the Fig.~1 caption now say so. Below the crossover
fixed overheads dominate, and the eight-way split at $N=40$ is $2.1\times$ \emph{slower} than a
single GPU, which is why the single-GPU sampler remains the right choice there.

Two rows carry caveats, both stated in the manuscript. $E=100.1\%$ at $N=54$ lies within the
run-to-run variability of these single-run timings and reads as ideal rather than superlinear
scaling, and $t_1$ at $N=40$ includes the one-off creation of the CUDA context, dropping to
$2.11$~s when re-measured after a warm-up.

\section*{Reviewer 2}

\point{R2.1}{The link to the script containing Simulated Bifurcation Machine (SBM) parameters
used for verification currently doesn't work. So, this section should contain the updated
link.}

\textbf{The data now ships with the article itself.} The
instances, the raw per-run results and the verification tables all live in the article
repository under \texttt{benchmarks/}:

\begin{center}
\small
\begin{tabular}{@{}ll@{}}
\texttt{benchmarks/instances/<N>.txt} & the exact instances used \\
\texttt{benchmarks/results/single-gpu/<N>.json} & per-run timings, energies, states \\
\texttt{benchmarks/results/distributed/<N>.json} & \\
\texttt{benchmarks/bf\_sbm\_verification\_*\_summary.csv} & per-size verification tables \\
\texttt{benchmarks/bf\_sbm\_verification\_*\_states.csv} & returned configurations \\
\texttt{benchmarks/sbm.py}, \texttt{verify\_sbm.py} & the comparison heuristic and its driver \\
\end{tabular}
\end{center}

The benchmark and plotting scripts sit alongside them and read from these paths, so Fig.~1 and
Table~3 regenerate from a single checkout with no external repository involved. Section~3 of
the revision is a guide to that directory.

On the SBM parameters specifically, the solver itself now ships with the article as
\texttt{benchmarks/sbm.py}, so there is no external script to point at and nothing to go
stale. Its configuration is also stated inline in the paragraph \emph{Verification against a
simulated-bifurcation solver}: \texttt{float64}, $2^{12}=4096$ parallel replicas, $3{,}000$
integration steps, seed 42, a global time step selected from $\{0.25,0.5,0.75,1.0,1.25\}$ by a
probe, coupling scale $c_0=0.5/(\sqrt{N}\,\mathrm{std}(J))$, and the lowest-energy replica
reported. Because all replicas are integrated simultaneously, one step is a single dense
$(N\!\times\!N)(N\!\times\!4096)$ product and the whole run costs
$\mathcal{O}(N^{2}R\,\times\,\mathrm{steps})$, independent of the exponential cost of
\emph{certifying} the optimum and a matter of seconds even at $N=60$, so a reader with no GPU
can reproduce the entire verification from this repository alone.

That also widened the claim. The heuristic is now compared against every stored single-GPU and
distributed sweep result underlying Fig.~1, all twenty runs at $N=38,\dots,60$ rather than the
nine of the submitted version, and on each of them it returns a configuration identical to the
certified ground state, with energies agreeing to $6\times10^{-14}$.

\begin{internal}
\internalhead
This answer requires the article repository to be \textbf{publicly reachable} at the URL given
by \texttt{\textbackslash datarepo} in the manuscript preamble (or archived with a DOI) at the
moment of resubmission. Verify before sending. This is the only outstanding item for R2.1:
\texttt{benchmarks/} now holds a single set of verification tables,
\texttt{bf\_sbm\_verification\_*}, produced by the shipped solver.
\end{internal}

\point{R2.2}{While the code achieves near ideal strong scaling for $N \ge 44$, corresponding
weak scaling tests are absent \dots\ This should help quantify the slowdown associated with the
controller merge step \dots\ which should ideally scale as the logarithm of processor count
\dots\ running this algorithm on $2^{16}$ processors could result in an algorithm that is
limited by this communication step.}

\textbf{We added the requested weak-scaling study and measured the local controller
components, while making the remaining network and large-allocation limitations explicit.}

The weak-scaling study follows the design the reviewer proposed. Table~4 reports
$N=50,51,52,53$ at $1$, $2$, $4$ and $8$ GPUs, with $k$ growing alongside the device count so
that every GPU always enumerates $2^{50}$ configurations. The wall-clock time stays between
$2112.5$~s and $2116.9$~s, flat to within $0.2\%$, with the merge contributing $0.007$~s
throughout. The same table reports strong scaling on the same layouts, both halves measured
twice on separate days.

The controller needs no large allocation to study, because the number of subproblems is set by
\texttt{num\_fixed\_vars} and not by the device count, so it can be driven on a CPU with the
partial results a worker would have returned. Table~5 reports that sweep to $2^{16}$
subproblems. The direct merge is linear, not logarithmic, as the reviewer anticipated: a
log--log fit gives $O(P^{0.99})$ over four decades. The merge arithmetic is nonetheless not the
limiting term, accounting for $2.1$~s of the $7.4$~s direct-merge total at $2^{16}$
subproblems, while dispatching and collecting $2^{16}$ \emph{empty} tasks already costs
$5.4$~s, more than the entire hierarchical merge. Task grouping is therefore listed as future
work, without claiming that scheduling alone sets the large-allocation ceiling.

Table~5 runs on one machine, so it excludes network transfer entirely, and both its text and
its caption present it as a lower bound on the controller cost at a given $P$. To probe what it
omits, the reproduction package adds \texttt{benchmarks/exp\_boundary\_cost.py}, which times
matched tasks that construct the same payload but only one of which returns it, under local,
remote and mixed placement. For batches of sixteen, the paired remote-minus-local effect was
$196.45$\,$\mu$s per task at $50{,}165$ bytes and $4.404$\,ms at $496{,}575$ bytes. That is the
aggregate ordinary Ray return path on two nodes, not an isolated network time, and we neither
fold it into Table~5 nor extrapolate it to thousands of raylets. The record ships as
\path{benchmarks/results/boundary_cost/2x1_production_20260816T122920_076038Z.json}.

That uncertainty is the architectural case for hierarchical reduction, which leaves the
controller receiving one object instead of $P$. We have implemented it. It overtakes the direct
merge only from $P=64$ upwards, since below that the concatenation is cheap enough that
performing it in remote tasks merely adds scheduling rounds, and the released sampler selects
between the two on that measured threshold.

The revision also states what the reported times include: the CUDA search for the single-GPU
sampler, and dispatch plus completion of every subproblem search for the distributed one, with
the final merge excluded and now reported separately in the sampler's metadata. Both exclusions
are polynomial in $N$ and negligible at the reported sizes. The trade-off in $k$ is explicit
too, since incrementing it at fixed $N$ halves the work per subproblem while doubling the
results to combine: a break-even point exists, and our measurement locates the merge-strategy
crossover but not the full solver's search-versus-controller break-even on a large allocation.

\point{R2.3}{While I could locate the code metadata table, I couldn't find the corresponding
software metadata table.}

\textbf{We have added the requested table.} The revision contains the \emph{Current executable software version} table
(S1--S7) as Table~2, next to the existing code metadata table.

\begin{itemize}[leftmargin=*]
\item \textbf{S2}: the PyPI distribution is a source archive only, so the ``executables'' are
compiled at installation time against the user's CUDA Toolkit, which the table states
explicitly.
\item \textbf{S5}: lists the CUDA Toolkit version used in continuous integration and the one
used for the reported measurements ($12.4$), together with \texttt{ray >= 2.9}, which is
required by the distributed sampler.
\end{itemize}

\begin{internal}
\internalhead
S6 (link to user manual) still points at the framework documentation and is contingent on the
publishing step under R3.7.
\end{internal}

\point{R2.4}{I couldn't locate \texttt{code/bf.py} in the directory containing the repository
(C2). This is presumably also available at [omni-bench], but the link doesn't work currently.
Also, in Section 2, it is stated that this benchmarking sweep is in C2, but I presume this is
not the case as other areas in the manuscript reference [omni-bench].}

\textbf{We thank the referee for spotting this, and the attribution is now consistent.}
Section~2 reads: the plugin sources and documentation are in the plugin repository (C2); the
benchmark and plotting scripts are in the article repository under \texttt{benchmarks/}; and
the instances, results and verification tables ship alongside them in the same directory. The
paths inside the scripts have been updated accordingly, so no script refers to an external
repository any more.

\point{R2.5}{The manuscript could also benefit from providing speedup with float32 numbers for
\texttt{num\_states = 1} by conducting the simulation at the default float64 precision. This
will help users \dots\ quantify the expected slowdown.}

\textbf{We measured the requested quantity.} Repeating the ground-state search at
$N=40,\dots,46$ in both precisions on one H100 gives, after excluding the $N=40$ warm-up
outlier, a \texttt{float64} penalty of only $2.0$--$2.3\%$ (ratios of $1.02$ at
$N=42,44,46$), with both precisions returning the same configuration at every size. The
enumeration is bound by memory traffic and integer bookkeeping rather than by floating-point
throughput, so on this hardware single precision buys working-set size, not speed, and
the \emph{Architecture} paragraph now says so.

Three code defects along the same path have been fixed. The SPIN$\,\to\,$BINARY recursion in
the distributed sampler dropped the requested \texttt{dtype}, so SPIN inputs always ran in
single precision; the command line's \texttt{--dtype} had no effect, because NumPy resolves
both \texttt{"float"} and \texttt{"double"} to \texttt{float64}, leaving the stabilized
single-precision path unreachable from the CLI; and the COO reader used by the benchmark
scripts dropped any coupling written in exponent notation, which affected one coupling in each
of the stored $N=46$, $58$ and $60$ instances. An exact parser now replaces that reader, and
the manuscript documents the resulting offset in the three affected records.

With the reader replaced, the precision picture is clean: the \texttt{float64} accumulator
stays within $10^{-13}$ of the from-scratch recomputation at every measured size, against
$1$--$4\times10^{-6}$ for stabilized \texttt{float32}, and the apparent \texttt{float64} drift
of $2.7\times10^{-5}$ at $N=46$ seen in an earlier version of this experiment was exactly that
dropped coupling rather than accumulation error. Both precisions return the same configuration
at every size, confirmed independently. Our guidance is unchanged: use the
default \texttt{float32} path and recompute the energy of the returned configuration in
\texttt{float64} on the host when a certified numerical value is needed. Extending the
stabilization to the (currently unstabilized) double-precision path is listed as future work.

\section*{Reviewer 3}

\point{R3.1}{$N$ is never defined, even though it plays a central role in interpreting the
results. It would be helpful to briefly reintroduce Omnisolver and its key parameters in the
introduction.}

\textbf{Both points are addressed.} $N$ is now defined at its first substantive use in the
abstract, again in the second paragraph of Section~1 and once more in the caption of Fig.~1,
as the number of binary variables of an instance (equivalently, the number of Ising spins), so
that exhaustive search enumerates all $2^{N}$ configurations.

The same paragraph introduces the sampler parameters that previously appeared in the code
listing without explanation: \texttt{num\_states}, \texttt{num\_fixed\_vars}$\,=k$,
\texttt{suffix\_size} (with the meaning that $2^{\texttt{suffix\_size}}$ configurations are
resident in the GPU working set at a time and the search sweeps
$2^{N-\texttt{suffix\_size}}$ such chunks), and the launch parameters \texttt{grid\_size},
\texttt{block\_size}, \texttt{num\_steps\_per\_kernel} and
\texttt{partial\_diff\_buffer\_depth}. Section~1 also opens with a short reminder of how
Omnisolver plugins are registered and what the framework provides.

\point{R3.2}{The omnisolver-bruteforce plugin is already mentioned in the original paper [1].
Therefore, statements such as ``this update introduces omnisolver-bruteforce'' \dots\ are
misleading. The paper should clearly state what is actually introduced by the software update
described here.}

\textbf{The reviewer is correct, and the contribution claim has been rewritten from scratch.}
The original publication lists the plugin, notes that it can be GPU-accelerated, reports a
benchmark of the GPU sampler on $1$, $2$, $4$ and $8$ NVIDIA A100 GPUs, and describes the
plugin as usable up to about $50$ variables. Neither ``first-class Omnisolver
plugin'' nor ``distributed multi-GPU execution'' is therefore presented as new any more:

\begin{itemize}[leftmargin=*]
\item The \textbf{title} changed to ``\dots\ -- a numerically stabilized, distributed GPU
exhaustive-search plugin'', which also removes the double colon.
\item The \textbf{abstract} now says the plugin, ``together with a preliminary multi-GPU
benchmark, was already announced in the original publication'' before listing what this update
adds.
\item A \textbf{new paragraph}, \emph{Relation to the original publication and to the
predecessor solver}, states the provenance explicitly (the plugin entry and the
1/2/4/8-A100 benchmark in ref.~[1], and the single-GPU CUDA solver of ref.~[2] underneath it),
and concludes that what the update contributes is not the existence of a GPU-accelerated or
multi-GPU exhaustive search, but its maturation into a dependable certification backend.
\item The \textbf{contribution list} was rewritten around what is verifiable in the released
code: a released, versioned, API-stable distributed sampler in place of the unreleased code
path behind the earlier benchmark; numerical stabilization of the \texttt{float32}
ground-state path; 64-bit-safe enumeration bookkeeping; certification reach extended
from $N\approx50$ to $N=60$, i.e.\ three orders of magnitude in enumerated configurations, with
an explicit strong-scaling analysis; and end-to-end verification of the returned optima.
\item The \textbf{conclusions} open with the same framing rather than with ``promotes \dots\ to
a first-class Omnisolver plugin''.
\end{itemize}

\point{R3.3}{To properly assess the scalability of the distributed sampler, experiments should
include more than two nodes \dots\ it would be interesting to investigate whether the
controller becomes a bottleneck as the number of nodes increases. For these reasons, the
scaling outlook also appears somewhat optimistic.}

\textbf{We agree on all three counts, and two of them we can now answer with measurements
rather than with argument.}

\emph{Does the controller become a bottleneck?} Not through its arithmetic. Table~5 measures
the local cost up to $2^{16}$ subproblems: $7.4$~s collecting directly, $3.8$~s reducing
hierarchically, against Ray's $5.4$~s task-scheduling floor for the same number of tasks, so
scheduling already exceeds the hierarchical merge before any network is involved. We claim no
more than that: as set out under R2.2, the measurement excludes the transfer of $P$ partial
results into one node and any contention at the head, which are the terms we expect to bind.
The revision keeps Table~5 as a local lower bound and leaves many-node transfer and contention
explicitly unmeasured.

\emph{Does the node layout matter to full-solver wall-clock time?} Not measurably. Holding the
GPU count fixed and adding a node boundary ($2\times1$ against $1\times2$, and
$2\times2$ against $1\times4$) changes
the wall-clock time by $-0.03\%$ and $-0.12\%$ respectively. The two-node configuration is
slightly faster in these runs, but the differences are below the experiment's resolution.

We measured the series twice to establish that resolution. A first pass gave $+1.47\%$ and
$+1.12\%$. Repeating it on a different day showed single-node configurations reproducing to
within $0.08\%$ across both sweeps ($0.02\%$ in the strong-scaling half) and two-node ones
moving by up to $1.5\%$, which places both figures inside the repeatability of two-node
timings. A single node boundary therefore changes full-solver wall-clock time by less than we
can resolve, with the $1.5\%$ set by that repeatability rather than by any systematic effect.
That is consistent with the absence of inter-worker communication during the search, and
localizes the variability to the two-node/second-host path, not to a particular
mechanism. The top-level E1/E2 result files contain the quoted second pass, and the first H100
pass is retained in the article repository history at commit \texttt{d965e54}.

\emph{More than two nodes} is the one question our hardware cannot settle: two nodes with four
H100s each on standard Ethernet is the entire cluster available to us. \emph{Limitations and
intended scope} states that, and extending the study past two node boundaries is now one of
the stated next steps in the conclusions, since a return-path effect or an effect below our
$1.5\%$ bound could still compound over many boundaries.

The scaling outlook itself is now $2^{14}$ rather than $2^{16}$ GPUs, as described under R1.1,
with $2^{16}$ retained only as an explicitly hypothetical bound far beyond the hardware
available for this study.

\point{R3.4}{QUBO instances are not evaluated in the paper, even though they are an important
target use case for the framework.}

\textbf{The measured path is in fact the QUBO path, which the revision now makes explicit.}
Both samplers convert a SPIN model to its BINARY equivalent, solve it as a QUBO on the GPU and
convert the result back, so the timings in Fig.~1 are QUBO timings. The only Ising-specific
cost is an $O(N^{2})$ host-side transformation outside the measured region. Loading the same
coefficient files once with each variable type gives a BINARY-to-SPIN wall-clock ratio of
$0.998$--$1.000$ over $N=42,\dots,50$. The energies differ, as they must, since the same
coefficients define different objectives under the two variable types. The runtimes do not.

\point{R3.5}{The paper should also specify which version of the CUDA Toolkit was used. \dots\
Which versions of the CUDA Toolkit are supported?}

\textbf{We have added the requested information.} The reported measurements were obtained with
\textbf{CUDA Toolkit 12.4}, and the revision-experiment drivers record the toolchain, driver,
GPU and package descriptors with their outputs (the original Figure~1 files predate that
field). On supported versions, the plugin is built and tested in continuous integration on a
$12.4$ and $12.5$ matrix under Python 3.10 and 3.11. C6 and S5 now state that range rather than
a single version, and S4 records the target architecture of the measurements (NVIDIA H100
96\,GB, compute capability \texttt{sm\_90}).

\point{R3.6}{The license listed in C3 does not exist in C2.}

\textbf{Corrected in the repository.} The Apache-2.0 license text is now present, declared in
the package metadata (\texttt{license} field) and included in the distribution. C1--C3 and
S1--S3 refer to the release that contains it.

\begin{internal}
\internalhead
The repository work is done (\texttt{LICENSE}, \texttt{pyproject.toml} metadata,
\texttt{MANIFEST.in}, SPDX headers in 5 files; \texttt{validate-pyproject} passes), but it is
\textbf{not released yet}. Before sending: (1)~confirm the copyright holder, since the headers
currently say ``The Omnisolver developers'', the alternative being the institutions;
(2)~merge \texttt{lp/add-docs} and \texttt{lp/better-merge}, tag \texttt{0.0.6} (the version
is derived by \texttt{setuptools\_scm}, so the tag is the release) and upload the sdist to
PyPI. C1, C2, S1, S2 and the contribution list already say \texttt{0.0.6}, so until the tag
exists and is on PyPI, the paragraph above and those table entries are a promise rather than
a fact.
\end{internal}

\point{R3.7}{The documentation linked in C7 is currently almost empty, contains several 404
errors, and is of little practical use in its current state.}

\textbf{We agree, and the documentation has been rewritten.} The former single page was a copy
of the \emph{framework's} landing page, describing Omnisolver in general and installing a
different plugin in its quickstart. The plugin now has its own front page and a user guide
covering installation, single- and multi-GPU usage, the meaning and selection of every sampler
parameter, the size limits, the precision and stabilization semantics and the reported timings.
The reference manual is generated from the docstrings, and the build is exercised in continuous
integration.

\begin{internal}
\internalhead
Done in the repository: \texttt{docs/index.md} rewritten; \texttt{docs/userguide.md} written;
\texttt{mkdocs.yml} repaired (dead \texttt{plugins.md} nav entry, missing
\texttt{stylesheets/extra.css}, wrong \texttt{repo\_url}, defunct polyfill.io, fragile
\texttt{with-pdf}); two configuration errors that broke the build on their own fixed (renderer
options nested under \texttt{docstring\_options}, and a Sphinx dependency stack declared for an
mkdocs project); \texttt{.readthedocs.yml} removed (it never built, and Read the Docs
\emph{cannot} build these docs because mkdocstrings imports the compiled CUDA extension),
replaced by a GitHub Pages workflow using \texttt{mike} on the CUDA-capable runner.
\texttt{mkdocs build --strict} succeeds. Before sending: run the workflow, put the published
URL into C7 and S6, and archive or redirect the stale Read the Docs project, since deleting the
config does not remove the site.
\end{internal}

\point{R3.8}{There are still a few minor typographical and grammatical errors throughout the
manuscript. A careful proofreading would help eliminate them.}

\textbf{The manuscript has been proofread throughout.} The corrections include the apposition
in ``certified reference, i.e.\ lowest-energy, solutions'', ``compute'' used as a noun, the
colloquial ``is clean already at $N=40$'', the typography of ``H100\,/\,96\,GB'', the
capitalization of ``Parallel Tempering'' and the double colon in the title (addressed with
R3.2). Spelling is now consistently American English.

\end{document}
